\section{Conclusion}\label{sec:conclusion}

This review has argued that verification, not capability, is now the binding constraint on autonomous agents that build software and do research. The same agent harness appears among the successes and the failures, running models generations apart. What separates them is who decides that the work is done: an executed test, a held-out metric or the running system, or the agent's own account and a reader of its prose. Two literatures that developed largely in parallel point to the same repairs. Software engineering offers executed checks in environments the builder cannot change, and the replication-crisis literature offers commitments made before results exist and independent re-execution.

Three kinds of finding would show the argument wrong. A budget-matched study in which self-graded agents matched independently graded ones on hidden checks of the running system would weaken the case for separation. Judges that read trajectories and matched execution-based checks in precision, across samples drawn by groups other than their builders, would weaken the case against internal signals. And if pass\textasciicircum k rose with capability until single runs ranked agents reliably, the demand for repeated measurement would fade. The current evidence runs the other way on all three: completion reports diverge from verified results \citep{F20}, verifiers that read rather than run approve much failing work \citep{T18}, and capability gains have brought only small improvements in reliability \citep{F2}.

Verification structures should also outlast any one model generation, for three reasons. First, the failures they address are structural. In a many-analysts study, neither expertise nor the peer-rated quality of the analyses readily explained why human teams diverged \citep{R9}, and agents from several model families gamed checkers within their reach \citep{V3,V17}. Second, the structures are indifferent to the model: a test the builder cannot edit, a decision rule committed in advance and an audited grader work whichever model builds. Third, as benchmarks saturate, as one repaired split already nearly has \citep{C12}, the question that remains is whether the instrument is right, which is a question about verification. Capability will keep rising; the structures that decide whether its output can be trusted are worth building now.
